\section{Control del resto al integrar}
Sea $X=m/T$ con $T$ de ley gamma de forma $m+1$ y tasa uno.
La identidad física es $m^m\per D'/m!=\E R_B(X)$.
Tomemos $I=[1-d,1+d]$ dentro del disco fijo $|y-1|<d_0$ de holomorfía
de los aproximantes, con $0<d<d_0$ y $(1+d)r_1r'<1$.
Para esta misma ventana la prueba de Divoux--Wang, sección 3.3, da
\[
 \E[|R_B(X)|;X\notin I]\le C e^{-c m+C\sqrt m}.
\]
En efecto su integral positiva se domina, después de un cambio de variable, por
la cola gamma desplazada por $W=\sum\sigma_i |Z_i|^2$, con peso $e^W$.
Los $|Z_i|^2$ son exponenciales de media uno y $\sigma_i$ los valores singulares de $B$.
En $W\le d'm$ la cola gamma cuesta $e^{-cm}$ para $d'$ pequeño según $I$.
En $W>d'm$ se usa $\E e^{(1+\gamma)W}\le e^{C\sqrt m}$ con $(1+\gamma)r'<1$.
Esto controla el valor absoluto del integrando incluso lejos del disco de Wick.
\paragraph{Función explícita.}
Para $k\le Q$ ponemos
\[
 H_k(u)=\sum_{i=0}^{2k+1}\mu_i(u)\frac{\partial_y^i S_k(1,u)}{i!},\qquad
 (1-iu)\mu_{i+1}=iu(2\mu_i+\mu_{i-1}),\quad \mu_0=1,\quad \mu_1=0 .
\]
Aquí $i$ es un índice entero. En $u=1/m$ los $\mu_i$ son los momentos de $X-1$
por integración por partes en la densidad de la gamma inversa.
Para $N|u|$ suficientemente pequeño, $t=N|u|/(1-N|u|)\le1/4$; por inducción
$|\mu_i(u)|\le(2\sqrt t)^i$ para $i\le N$. En el centro físico esto da
$\E|X-1|^{2h}\le (Ch/m)^h$ para $2h\le c_1m$, usando la potencia par.
Tomamos $4Q+4<c_1m$ y reducimos el radio $|u|\le c_0/k$ según las constantes.
Por Cauchy en un radio fijo $r_0$ con $d<r_0<d_0$, los coeficientes de Taylor de
$S_k$ están acotados por $Cr_0^{-i}$.
En $I$ el error al integrar su Taylor hasta $2k+1$ cuesta
$C r_0^{-2k-2}\E|X-1|^{2k+2}/(1-d/r_0)$.
En $I^c$ cada monomio de grado $i\le2k+1$ satisface
$|X-1|^i\le d^{i-2k-2}|X-1|^{2k+2}$; la suma geométrica usa $d/r_0<1$.
Así sustituir $\E[S_k(X,1/m);I]$ por $H_k(1/m)$ cuesta
$(C(k+1)/m)^{k+1}$. Para $k=Q$ es $O(e^{-c'j})$ reduciendo $\epsilon$.
La parte 03 y la cola ponderada dan entonces el resultado físico con $k=Q$.
Para bajar a $q$, en $I$ se tiene
$\sup|S_Q-S_q|\le C(C(q+1)/m)^{q+1}$, pues los mayorantes $(Cl/m)^l$
tienen razón sucesiva $\le Ce(l+1)/m\le1/2$. Se aplica el mismo Taylor a $q$.
\paragraph{Paso al logaritmo.}
También para datos complejos y $|u|\le c_0/k$,
$\sum_{l=1}^k(Cl|u|)^l=O(|u|)$: tras extraer $|u|$ la serie está dominada por
$C\sum_{l\ge1}l(Cc_0)^{l-1}$. En la cota de $\mu_i$ anterior tomamos $N=i$
para cada $i\ge2$. Tras extraer $|u|$ y multiplicar por $r_0^{-i}$ se suma una
serie acotada por $C\sum_{i\ge2}i\theta^{i-2}$ con $\theta<1$ fijado mediante $c_0$.
Como $\mu_1=0$, resulta uniformemente $H_k(u)=d(1)+O(|u|)$.
La parte 02 separa $d(1)$ de cero y fija la rama compleja; en los datos físicos el
permanente normalizado es además $\ge1$ por van der Waerden.
Los errores absolutos anteriores pasan a errores de logaritmos con constantes uniformes.
El caso $B=0$ satisface la fórmula con $R_B=d=1$.
